\section{Trabalhos Relacionados}

\subsection{Bases e tarefas em citologia cervical}

Bases como Herlev e SIPaKMeD são frequentes em estudos de classificação
cervical \cite{jiang2023systematic,fang2024systematic}. A SIPaKMeD reúne
células extraídas de esfregaços de Papanicolau para classificação por atributos
e imagens \cite{plissiti2018sipakmed}. Embora úteis para comparação, bases de
células isoladas não reproduzem integralmente sobreposição, inflamação,
variação de fundo e outros fatores encontrados em exames convencionais.

A CRIC Cervix foi construída a partir de 400 imagens de citologia convencional
e contém 11.534 células classificadas por citopatologistas
\cite{rezende2021cric}. A coleção inclui NILM, ASC-US, LSIL, ASC-H, HSIL e
SCC, preservando desafios de lâminas reais, como células sobrepostas e
elementos inflamatórios. Por isso, é útil para estudar variação entre imagens,
desde que a dependência entre células da mesma imagem seja considerada.
Essa característica diferencia a CRIC de bases compostas por células isoladas:
o modelo pode explorar tanto atributos celulares quanto padrões de aquisição,
coloração e fundo compartilhados por uma mesma imagem. Assim, o particionamento
por imagem é parte central do desenho experimental.

\subsection{Aprendizado profundo para classificação cervical}

Estudos recentes reportam desempenho elevado com redes profundas em Herlev,
SIPaKMeD e outras bases de citologia cervical. Kaur et al. compararam modelos
de transferência e reportaram 95\% de acurácia com ResNet50 em Herlev binário
\cite{kaur2025comparison}; Mehedi et al. obtiveram 98,53\% de acurácia com
CCanNet em SIPaKMeD de cinco classes \cite{mehedi2024lightweight}; e Lima et
al. reportaram 98,35\% em Herlev binário e 99,73\% em SIPaKMeD binário com
\emph{ensembles} \cite{lima2025ensemble}. Especificamente com ConvNeXt, Tang et
al.\ reportaram alta precisão na classificação de lesões precursoras cervicais,
evidenciando o potencial dessa família em citologia \cite{tang2023convnext}.
Esses resultados são relevantes como
contexto, mas não são diretamente comparáveis ao presente estudo por diferenças
de base, preparação, número de classes, métrica principal e unidade de
particionamento. Assim, mais do que confrontar percentuais absolutos, a
comparação com a literatura deve considerar se a separação entre treino e teste
ocorre no nível de célula, imagem, lâmina ou paciente.

Em citologia convencional, Rodríguez et al. reportaram sensibilidade 0,688 e
especificidade 0,762 em esfregaços de Papanicolau, indicando maior dificuldade
que em bases de células isoladas ou citologia em meio líquido
\cite{rodriguez2024automated}. Este trabalho se diferencia pelo protocolo:
CRIC, separação por imagem, validação cruzada agrupada, ajuste fino da
ConvNeXt-Tiny em cinco folds e apresentação explícita da variabilidade entre
partições.

Além do desempenho, o relato de estudos de IA em saúde deve permitir
interpretação crítica do risco de viés, da independência dos dados e da
aplicabilidade do modelo \cite{collins2024tripod,tejani2024claim}. Essa
preocupação é especialmente relevante em citologia baseada em recortes, pois
várias amostras podem vir do mesmo campo microscópico.
